% ============================================================
% SECCIÓN 7: Testing y Validación
% ============================================================
\chapter{Testing y Validación}
\label{sec:testing}

% ─────────────────────────────────────────────────────────────
\section{Estrategia de Testing por Capa}

El sistema Demeter está formado por cuatro capas con tecnologías y restricciones de
ejecución radicalmente distintas: firmware C\texttt{++} sobre ESP32, servicio Python
asíncrono en la Raspberry Pi, backend FastAPI en GCP y SDK científico empaquetado para
el usuario final. Aplicar una única política de testing a todas ellas no sería
realista; en su lugar, cada capa adopta el enfoque que mejor se ajusta a sus
dependencias y a la madurez de los componentes que ejercita.

\begin{table}[H]
    \centering
    \caption{Enfoque de testing por capa del sistema}
    \begin{tabularx}{\textwidth}{lllX}
        \toprule
        \textbf{Capa}            & \textbf{Framework}      & \textbf{Énfasis}     & \textbf{Justificación} \\
        \midrule
        Firmware (ESP32)         & Unity (PlatformIO)      & Unitario             & Hardware caro de simular: se aíslan controladores y motor de protocolo con mocks Arduino. \\
        Raspberry Pi             & pytest, asyncio         & Mixto                & Lógica de despacho concurrente: se mockean transporte UART y WebSocket para validar flujos sin red ni puerto serie. \\
        Backend (FastAPI)        & pytest, TestClient      & Integración           & La autenticación JWT, el RBAC y la cola RQ son emergentes: se ejercitan extremo a extremo contra una base SQLite en memoria. \\
        SDK Científico           & pytest                  & Unitario              & Funciones agronómicas puras (VPD, entalpía, GDD): se validan con vectores numéricos de referencia, sin red. \\
        \bottomrule
    \end{tabularx}
    \label{tab:test_strategy}
\end{table}

\begin{designnote}
\textbf{Pirámide de testing pragmática.} El proyecto privilegia la \textit{base} de la
pirámide — tests unitarios rápidos, deterministas y locales — y un nivel intermedio de
integración acotado a fixtures en memoria. No se invierte en E2E automatizados (frontend
clic-a-clic, hardware físico continuo) porque su mantenimiento no se amortiza en un
sistema de un solo desarrollador. La validación E2E se delega a \textit{spikes}
manuales documentados (Sección~\ref{sec:testing_field}) y a la operación real del
sistema durante los 30 días de estabilización.
\end{designnote}

% ─────────────────────────────────────────────────────────────
\section{Tests Unitarios}
\label{sec:testing_unit}

La capa unitaria cubre los componentes cuyo contrato es internamente verificable: un
controlador GPIO al que se le piden estados, un parser binario al que se le suministran
tramas, un cálculo agronómico al que se le pasan vectores. Todos estos tests se
ejecutan en menos de un segundo y no requieren hardware ni red.

\begin{table}[H]
    \centering
    \caption{Inventario de tests unitarios (verificado el 2026-05-07)}
    \begin{tabularx}{\textwidth}{lXrl}
        \toprule
        \textbf{Capa}           & \textbf{Archivo}                      & \textbf{\#Tests} & \textbf{Cubre} \\
        \midrule
        Firmware                & \texttt{test\_gpio\_controller.cpp}   & 3   & Lectura/escritura de pines, pull-up, modo \\
        Firmware                & \texttt{test\_protocol\_engine.cpp}   & 25  & Parseo, CRC, cabecera, callbacks (MockComms) \\
        Firmware                & \texttt{test\_system\_manager.cpp}    & 4   & Inicialización, transición de estados \\
        Raspberry Pi            & \texttt{test\_device\_manager.py}     & 4   & Alta, baja y consulta de nodos \\
        Raspberry Pi            & \texttt{test\_telemetry\_cache.py}    & 23  & Caché en memoria, expiración, agregación \\
        Raspberry Pi            & \texttt{test\_ws\_client.py}          & 3   & Reconexión, serialización de frames \\
        SDK                     & \texttt{test\_science.py}             & 23  & VPD, entalpía, presión de vapor, GDD \\
        SDK                     & \texttt{test\_transform.py}           & 15  & Pipeline ETL, resample, normalización \\
        SDK                     & \texttt{test\_client\_offline.py}     & 7   & Caché Parquet sin red \\
        SDK                     & \texttt{test\_export.py}              & 7   & Exportación a CSV y Excel \\
        \midrule
        \textbf{Total}          &                                       & \textbf{114} & \\
        \bottomrule
    \end{tabularx}
    \label{tab:tests_unit}
\end{table}

\textbf{Detalles destacables por capa:}

\begin{itemize}[leftmargin=*, label=--]
    \item \textbf{Firmware}. El archivo \texttt{test\_protocol\_engine.cpp} concentra
          25 casos sobre el parseo del protocolo binario Demeter: validación de la
          cabecera (Sección~\ref{sec:protocolo}), CRC, manejo de tramas truncadas,
          callbacks y rutas de error. Es el componente con mayor densidad de tests
          porque es también el más expuesto a entradas externas.
    \item \textbf{Raspberry Pi}. \texttt{test\_telemetry\_cache.py} (23 tests) ejercita
          la caché que evita golpear al backend en cada lectura. Cubre expiración por
          TTL, agregación por nodo y comportamiento bajo carga concurrente con
          \texttt{asyncio}.
    \item \textbf{SDK}. Las funciones científicas son las que más tests soportan
          porque su corrección numérica es auditable: cada función agronómica se
          contrasta contra al menos un caso de referencia conocido (e.g.\ tabla de
          presión de vapor de saturación, ejemplos del manual de agroclimatología).
\end{itemize}

% ─────────────────────────────────────────────────────────────
\section{Tests de Integración}
\label{sec:testing_integration}

La capa de integración ejercita la \textit{colaboración} entre componentes: una
petición HTTP que atraviesa autenticación, base de datos y respuesta; un comando que
viaja del despachador al transporte UART; una tarea que se encola y se procesa.
Ningún test de integración requiere infraestructura real: PostgreSQL se sustituye por
SQLite en memoria, Redis se mockea y el transporte UART se reemplaza por un
\texttt{AsyncMock}.

\begin{table}[H]
    \centering
    \caption{Inventario de tests de integración}
    \begin{tabularx}{\textwidth}{lXrl}
        \toprule
        \textbf{Capa}           & \textbf{Archivo}                       & \textbf{\#Tests} & \textbf{Ejercita} \\
        \midrule
        Raspberry Pi            & \texttt{test\_command\_dispatcher.py}  & 3  & Despachador + UART + WebSocket \\
        Raspberry Pi            & \texttt{test\_uart\_processor.py}      & 4  & Transporte UART + handshake (Ping/Ack) \\
        Backend                 & \texttt{test\_auth.py}                 & 6  & Login, refresh token, RBAC \\
        Backend                 & \texttt{test\_command.py}              & 4  & RBAC + JWT + envío de comando \\
        Backend                 & \texttt{test\_analysis.py}             & 3  & Encolado en RQ, procesamiento \\
        Backend                 & \texttt{test\_lims.py}                 & 3  & Endpoints LIMS sobre DB en memoria \\
        Firmware                & \texttt{test\_integration.cpp}         & 0$^{\dagger}$ & Bucle UART por loopback (infraestructura) \\
        \midrule
        \textbf{Total}          &                                        & \textbf{23} & \\
        \bottomrule
    \end{tabularx}
    \label{tab:tests_integration}
\end{table}

\noindent$^{\dagger}$ Infraestructura de loopback presente pero sin casos
\texttt{TEST\_CASE} activos: queda como base para los tests Hardware-In-the-Loop
previstos tras el spike de bombas (Sección~\ref{sec:testing_field}).

\textbf{Patrón compartido.} Todos los tests de integración del backend usan un
fixture \texttt{conftest.py} que arranca una aplicación FastAPI completa contra
una base SQLite en memoria, crea usuarios con distintos roles y emite tokens JWT
para ejercitar los endpoints. Los tests de la Raspberry Pi siguen un patrón análogo
con \texttt{AsyncMock} para los componentes asíncronos (Sección~\ref{sec:raspberry}).

\begin{warningbox}
\textbf{Limitación conocida.} Ningún test de integración cubre actualmente el flujo
completo \textit{usuario $\rightarrow$ backend $\rightarrow$ Raspberry Pi $\rightarrow$
ESP32 $\rightarrow$ actuador físico}. Esa validación se hizo manualmente sobre banco,
comprobando que el comando llega al nodo y conmuta el pin; nunca llegó a ejercitarse
contra un actuador hidráulico real. Automatizarla exigiría un \textit{harness} con un
ESP32 dedicado conectado a la suite de tests, un coste de mantenimiento que el
proyecto no asumió.
\end{warningbox}

% ─────────────────────────────────────────────────────────────
\section{Pruebas de Campo y Validación Hardware}
\label{sec:testing_field}

Las pruebas de campo cubren lo que el banco de tests automatizados no puede: el
comportamiento de los componentes electromecánicos reales bajo las condiciones de
operación. Estas pruebas se documentan como \textit{spikes} con objetivos, hipótesis
verificables y criterios de aceptación previos a la ejecución.

\textbf{SPIKE-001 — Validación de la bomba de 5\,V: planificado y no ejecutado.} El
primer y único \textit{spike} del ciclo se escribió para confirmar cinco hipótesis
sobre una bomba mini sumergible antes de comprometer el resto del sistema de riego:

\begin{itemize}[leftmargin=*, label=--]
    \item \textbf{H1}: la bomba entrega un caudal $\geq$~50~mL/min a 5\,V.
    \item \textbf{H2}: el consumo a régimen no excede 350\,mA, dentro del margen del módulo relé.
    \item \textbf{H3}: el ESP32 actuador puede activar/desactivar la bomba a través del relé sin glitches.
    \item \textbf{H4}: el tiempo de riego típico para una maceta (~50~mL) es inferior a 90\,s.
    \item \textbf{H5}: la temperatura del cuerpo de la bomba se mantiene estable tras 6 ciclos consecutivos.
\end{itemize}

El \textit{spike} quedó estructurado en tres sesiones de tres horas --- montaje,
caracterización y ciclos repetidos --- con una tabla de resultados que debía decidir
la continuidad: continuar, pivotar a otra bomba o abortar el rediseño físico. El
procedimiento y la plantilla residen en
\texttt{docs/spikes/SPIKE-001-validacion-bomba-5v.md}.

\textbf{Ninguna de esas tres sesiones se llegó a celebrar.} El \textit{spike} está
completo como documento y vacío como experimento: sus cinco hipótesis siguen sin
verificar, y la tabla de resultados sin rellenar. Peor aún, quedó desalineado del
hardware que finalmente se adquirió: está redactado para una bomba de 5\,V y la
instalación comprada es de 12\,V autocebada, de modo que sus umbrales de caudal y
consumo ya no son aplicables. Cualquier continuación tendría que reescribirlo antes
de ejecutarlo.

Se conserva aquí, sin maquillarlo, porque documenta el punto exacto en el que el
proyecto dejó de avanzar sobre el mundo físico, y porque un \textit{spike} escrito y
no ejecutado es una lección de proceso más útil que su ausencia.

\textbf{Validación E2E continua: no realizada.} El criterio de \texttt{DONE} previsto
para la versión 1.0 era un periodo de \textbf{30 días de operación continua sin
intervención manual}, durante el cual se anotarían caídas del servicio, consumo
energético, tasa de pérdida de tramas y desviación temporal del secuenciador. Ese
periodo nunca se inició. \textbf{El sistema acumula cero días de operación
desatendida}, y en consecuencia ninguna de esas cuatro series existe.

\textbf{Pruebas de campo previas (firmware).} El firmware ya ha sido sometido a
pruebas de Deep Sleep en banco que validan el ciclo energético documentado en el
Anexo (Figura~\ref{fig:anx_ciclo_energetico}): fase activa de $\sim$1{,}5\,s y
ventana de Deep Sleep de $\sim$3596\,s. Los detalles de medición están en
\texttt{docs/deep\_sleep\_test.md}.

% ─────────────────────────────────────────────────────────────
\section{Métricas Medidas}
\label{sec:testing_metrics}

Esta sección recoge únicamente métricas \textbf{efectivamente medidas} sobre el
sistema. Las métricas para las que aún no existe instrumentación se marcan como
pendientes para evitar inducir confianza injustificada en los datos.

\begin{table}[H]
    \centering
    \caption{Métricas observadas y estado de medición}
    \begin{tabularx}{\textwidth}{Xlll}
        \toprule
        \textbf{Métrica}                                  & \textbf{Valor}            & \textbf{Fuente}                                & \textbf{Estado} \\
        \midrule
        Tiempo de fase activa nodo sensor                  & $\sim$1{,}5\,s            & Banco con osciloscopio                          & Medido \\
        Periodo Deep Sleep nodo sensor                     & $\sim$3596\,s             & Configuración firmware + medida real            & Medido \\
        Consumo nodo en Deep Sleep                         & $\mu$A (orden de magnitud)& Datasheet ESP32 + observación                  & Medido \\
        Latencia comando usuario\,$\rightarrow$\,actuador  & ---                       & Sin instrumentación                            & Pendiente (spike) \\
        Throughput de tramas Gateway\,$\rightarrow$\,RPi   & ---                       & Sin instrumentación                            & Pendiente \\
        Tasa de pérdida de tramas                           & ---                       & Sin instrumentación                            & Pendiente \\
        Cobertura de código (todas las capas)              & No formalizada            & \texttt{pytest-cov} declarado en SDK, sin reporte & Pendiente \\
        \bottomrule
    \end{tabularx}
    \label{tab:metrics}
\end{table}

\textbf{Las métricas pendientes lo siguen estando.} El plan de instrumentación
previsto --- marcas de tiempo en frontend y ESP32 para la latencia extremo a extremo,
contadores de tramas recibidas y descartadas en el servicio de borde, y ejecución de
\texttt{pytest --cov} sobre backend, Raspberry Pi y SDK --- \textbf{no llegó a
implementarse en ninguno de sus tres puntos}.

La tabla anterior queda, por tanto, como el estado definitivo de la caracterización
del sistema: tres métricas medidas, todas ellas del ciclo energético del nodo, y
cuatro sin instrumentar. Es poco, y decirlo importa: sin línea base medida no hay
forma de demostrar que el sistema cumple sus requisitos ni de detectar una regresión.
Cualquier continuación de este trabajo debería empezar por aquí, antes de añadir una
sola función nueva.

% ─────────────────────────────────────────────────────────────
\section{Resumen del Estado de Validación}

A fecha de redacción de este capítulo, el banco de tests automatizado del proyecto
suma \textbf{140 tests verificables} distribuidos en \textbf{20 archivos} sobre las
cuatro capas del sistema y la interfaz web:

\begin{table}[H]
    \centering
    \caption{Resumen agregado de tests por capa}
    \begin{tabular}{lrrrr}
        \toprule
        \textbf{Capa}      & \textbf{Archivos} & \textbf{Tests} & \textbf{Unitarios} & \textbf{Integración} \\
        \midrule
        Firmware           & 5                 & 32             & 32                 & 0$^{\dagger}$ \\
        Raspberry Pi       & 5                 & 37             & 30                 & 7 \\
        Backend            & 4                 & 16             & 0                  & 16 \\
        SDK                & 4                 & 52             & 52                 & 0 \\
        Frontend           & 2                 & 3              & 3                  & 0 \\
        \midrule
        \textbf{Total}     & \textbf{20}       & \textbf{140}   & \textbf{117}       & \textbf{23} \\
        \bottomrule
    \end{tabular}
    \label{tab:tests_summary}
\end{table}

\noindent$^{\dagger}$ Infraestructura de loopback presente, sin casos
\texttt{TEST\_CASE} activos.

\begin{designnote}
\textbf{Honestidad metodológica.} El propósito de este capítulo no es exhibir un
banco de pruebas exhaustivo, sino documentar con precisión qué se ha verificado,
cómo se ha verificado y qué queda fuera. La cifra de 140 tests es modesta a escala
industrial; en cambio, cubre los componentes de mayor riesgo del sistema --- parser
binario, despachador asíncrono, autenticación --- que es donde un fallo silencioso
resulta más caro.

Lo que no puede decirse es que ese banco se complementara, como estaba previsto, con
validación manual estructurada y operación continua: el \textit{spike} de hardware no
se ejecutó y los 30 días de operación desatendida nunca empezaron. La versión 1.0
se entrega, por tanto, \textbf{verificada en software y sin validar en campo}. Es una
distinción que este documento prefiere dejar escrita antes que difuminar.
\end{designnote}
